% 14 · Spectral representations — reading the whole spectrum
\section{SPECTRA INSTEAD OF ELEMENT RATIOS}
\label{sec:spectral}

\deckslides{slide 47, the masked-autoencoder block and its four steps.}

Everything in \S\ref{sec:benchmark} treats a star as a 16-dimensional vector of
abundances. That vector is a lossy summary of the observation it came from: a
single APOGEE spectrum has 8\,575 wavelength pixels, and the abundance
pipeline that turns those pixels into 16 numbers assumes a one-dimensional
model atmosphere, a set of atomic and molecular line data, and a fitting
procedure that was tuned on the Sun. This chapter asks the obvious question ---
if the spectra are the data, why cluster the summary? --- and reports what
happens when we do.

\begin{marginnote}[]
\entry{Masked autoencoder}{A self-supervised model that hides part of its
input, encodes the visible part, and is trained to reconstruct the hidden
part. It learns the structure of the data without any labels. In 1-D spectra,
\emph{how} you mask decides what it can learn.}
\entry{Self-supervised}{Trained with a loss defined by the data itself
(reconstruction) rather than by external labels. The property we care about
here is negative: a self-supervised latent never saw an element ratio or a
cluster membership, so it cannot be circular.}
\end{marginnote}

\subsection{The baseline: PCA of the spectra}

The honest first experiment is principal component analysis. \citet{GarciaDias:19}
applied eight clustering algorithms to the 13 APOGEE abundances of each star
under four dimensionality-reduction strategies, and this chapter asks the same
question one level down: which representation of the observation keeps the
cluster information. Reading spectra with
learned models has a longer history in this field --- the Cannon builds a
data-driven spectral model from labels \citep{Ness:15}, and convolutional
networks do the same job without the labels' linearity assumption
\citep{Ting:19} --- and the question this chapter asks is which of those
representations keeps the cluster information. Keep 64 components and
the classification results of \S\ref{sec:benchmark} already improve on the
abundances; keep 256 and they get slightly worse. PCA is the right baseline and
the wrong tool for one reason that is worth internalising: it maximises
\emph{variance}, and in a stellar spectrum the variance is dominated by things
that have nothing to do with chemistry --- the temperature-sensitive continuum
shape, gravity-sensitive line profiles, interstellar reddening, and the
signal-to-noise of the observation. A principal component that encodes
$T_{\rm eff}$ is useful for almost every other problem in stellar
spectroscopy and actively unhelpful here, because a hot field star and a hot
cluster member will share it.

\subsection{The model: a masked spectral autoencoder}

Our proposal \citep{He:22} is to train a self-supervised latent that has to
predict the spectrum's hidden parts:

\begin{enumerate}
\item Take the continuum-normalised spectrum, 8\,575 pixels.
\item Mask 50\% of the wavelength axis, in \emph{contiguous blocks} of 200
  pixels --- roughly a spectral line region at APOGEE's resolution.
\item Encode the visible pixels with five convolutional blocks, each halving
  the wavelength axis, with channel widths from 1\,024 down to 64, ending in a
  global average pool and a 256-dimensional linear bottleneck.
\item Reconstruct the masked pixels with a mirrored decoder --- and compute the
  loss \emph{only where the input was hidden}.
\item Discard the decoder. Keep the encoder's 256-dimensional output and feed
  it to the clustering pipeline exactly as the abundances were.
\end{enumerate}

The model has 8.6 million parameters and trains in under an hour on one GPU on
39\,945 field spectra plus 994 member spectra, with early stopping. Nothing in
the loss function mentions abundances, clusters, kinematics, or age.

\begin{figure}[htbp]\centering
  \fig[0.60]{mae_arch.png}
  \caption{The architecture: five strided convolutional blocks compress the
  8\,575-pixel spectrum by a factor of 32 while widening the channel count, a
  global average pool removes the wavelength axis, and a linear layer produces
  the 256-dimensional latent. The decoder is a mirror of the encoder and is
  discarded at inference.}
  \label{fig:maearch}
\end{figure}

\subsection{Why the masking is contiguous}

This is the one design decision in the chapter that a reader should be able to
justify from first principles, and it is where a naive implementation fails.

Mask single pixels at random and a convolutional network does not have to
learn chemistry: the masked pixel's immediate neighbours are visible, and the
spectrum is smooth on the scale of a few pixels, so copying from the left and
right solves the task. The reconstruction loss falls, the latent looks
respectable, and the model has learned to interpolate. Mask \emph{contiguous
200-pixel blocks} and the visible pixels adjacent to the hole are far enough
away, in wavelength terms, to be in different lines; the model must use the
global shape of the spectrum --- which line regions are present, and how deep
--- to fill in the gap. That is where element ratios live.

\begin{figure}[htbp]\centering
  \begin{tabular}{@{}c@{\hskip6pt}c@{}}
    \fig[0.45]{mae_why_blocks.png} &
    \fig[0.45]{mae_step1_mask.png}
  \end{tabular}
  \caption{Left: why the masking is contiguous --- a single-pixel hole is
  filled by interpolation from its immediate neighbours, while a 200-pixel
  block removes a whole line region and forces the model to use context.
  Right: the actual mask used (top) and the residual the model must predict
  (bottom), with the masked region shaded.}
  \label{fig:maemask}
\end{figure}

\begin{figure}[htbp]\centering
  \begin{tabular}{@{}c@{\hskip6pt}c@{}}
    \fig[0.45]{mae_step2_encode.png} &
    \fig[0.45]{mae_step3_recon.png}
  \end{tabular}
  \caption{Encoding and reconstruction. Left: only the visible 50\% of the
  spectrum enters the encoder. Right: the decoder's reconstruction of the
  hidden half, against the truth, with the loss computed only on the hidden
  pixels. The model is never rewarded for the parts it can see --- which is
  what makes its latent a description of the spectrum rather than a copy of
  it.}
  \label{fig:maeenc}
\end{figure}

\begin{figure}[htbp]\centering
  \fig[0.55]{mae_step4_latent.png}
  \caption{The 256-dimensional latent, projected. The spectral latent carries
  the abundance structure of the members because it had to reconstruct the
  lines those abundances are measured from --- not because anyone told it she
  was looking at C, N, O or Fe.}
  \label{fig:maelatent}
\end{figure}

\subsection{What it achieves}

On cluster-only separation it is the strongest signal we have
(\textbf{Table~\ref{tab:honest}}): homogeneity $0.740 \pm 0.002$ (t-SNE),
$0.760 \pm 0.011$ (UMAP), $0.714 \pm 0.036$ (EVoC) on the 982-star population,
against abundances $0.560$/$0.578$/$0.420$. In the single-run baseline it
reaches $0.789$/$0.874$/$0.770$ on the full 1\,002-row matrix, and
concatenating it with the abundances gives $0.879$ --- the highest cluster-only
score measured anywhere in this project. It also has the most important
non-empirical property: because it never saw an abundance label, it cannot be
circular in the way a supervised latent is.

On field retrieval it is a different story, and the difference is the most
useful result in this chapter:

\begin{itemize}
\item \textbf{Facing the field, the masked latent does not beat PCA-64.} In
  \textbf{Table~\ref{tab:field}} the two are
  $0.418$/$0.619$ and $0.417$/$0.631$: a tie to within the third decimal.
  Two very different representations of the same spectra --- one unsupervised
  and linear, one self-supervised and deep --- land on the same operating
  point.
\item \textbf{The tie is the finding.} If a 64-dimensional linear projection
  reaches the same field-retrieval precision as a 8.6-million-parameter
  convolutional latent, then what is limiting the precision is not the
  representation capacity of the model. It is the information in the spectra
  themselves, or the way our metric uses it.
\item \textbf{EVoC disagrees, and we report both.} On the same cluster-only
  task with EVoC as the clusterer, the supervised CNN reaches
  $0.80 \pm 0.03$ on the four-cluster sample against the masked latent's
  $0.70 \pm 0.08$ --- the supervised model wins, and the two error bars
  overlap. We described the self-supervised result as a
  t-SNE/UMAP result, not as a universal one, and withdrew the stronger
  sentence.
\end{itemize}

\begin{figure}[htbp]\centering
  \fig[0.72]{paired_control.png}
  \caption{The paired control that ended the ``two times purer'' claim, which
  is the reason this figure is in the workbook. Each star is embedded through
  both pipelines; the panel shows the distance from a star to \emph{itself}
  across data products against the distance to a different random star within
  one product. The mean is indicated. When a star is further from its own
  second embedding than from a stranger, the latent is encoding something
  about the data product rather than about the star.}
  \label{fig:paired}
\end{figure}

\subsection{The error we made, in full}
\label{sec:provenance}

A masked autoencoder trained on a mixed set of spectra will learn whatever
distinguishes them, and our first training set was mixed in a way we did not
notice. The field sample was 100\% raw \texttt{apStar} spectra, while 91\% of
the member sample was continuum-normalised \texttt{aspcapStar} output, with
median flux levels differing by a factor of $\sim\!6\,000$. Along that boundary
the labels themselves are the difference.

The consequences, all measured:

\begin{itemize}
\item A linear probe recovers the data product from the latent at
  \textbf{AUC 0.999} on members alone, using 223 of the 256 dimensions
  ($p < 10^{-6}$).
\item A one-bit flag (``which product is this?'') separates members from field
  stars at \textbf{AUC 0.957} --- essentially the whole latent score, with no
  chemistry involved.
\item The paired control of \textbf{Figure~\ref{fig:paired}} is decisive: for
  253 stars embedded through both pipelines, a star sits $1.70\times$ further
  from its own other-product embedding than from a random different star
  within one product (cosine distance $0.389$, with $75\%$ of the offset in
  one shared direction --- the signature of a single global nuisance axis).
\end{itemize}

We first diagnosed this as a data-release effect (DR17 members, DR19 field).
That was wrong, and the check that disproved it is worth repeating: DR19
reanalyses and includes DR17 --- 717\,689 rows carry
\texttt{release=`dr17'} --- and DR19 serves a uniform \texttt{mwmStar}
spectrum for \textbf{738 of 738} members in the backfilled list. The mismatch
was between two products on the same sky, at the same epoch, from the same
spectrograph.

So the field-retrieval row for the mixed latent is not interpretable, and the
``twice as pure'' precision claim that depended on it is withdrawn. What
survives is narrower and solid: the cluster-only homogeneity, which is computed
between members and therefore within one product, holds on a single-product
subsample ($0.612$/$0.686$/$0.575$ against the abundances' $0.484$/$0.550$/
$0.488$ on the same stars). The fix is not a caveat but a re-run: download
uniform \texttt{mwmStar} spectra for every member and field star, re-embed, and
repeat the benchmark. That run is staged in the repository --- 736
\texttt{mwmStar} spectra are already downloaded and waiting --- and its results
belong in the next version of this workbook rather than in this one.

\begin{issues}[EXERCISES]
\begin{enumerate}
\item Train the same autoencoder twice, once with single-pixel random masking
  and once with contiguous 200-pixel blocks, and compare the reconstruction
  loss on a held-out set. Then compare the two latents on the cluster-only task.
  Which model wins on reconstruction, and which on chemistry? Explain the
  disagreement.
\item Reproduce the linear probe of \S\ref{sec:provenance}: fit a logistic
  regression from a 256-dimensional latent to a binary flag identifying the
  data product, and report the AUC on held-out member stars. Now repeat it for
  a latent trained on a single product. What is the AUC you would consider
  acceptable before publishing a field-retrieval number?
\item The paired control uses 253 stars with embeddings through both pipelines.
  Design an equivalent control for a study that only has one product, and
  state what it can and cannot detect.
\item The masked latent and PCA-64 tie on field retrieval. Working from the
  two latents, characterise what each one encodes: fit a linear model from
  each latent to $T_{\rm eff}$, $\log g$ and [Fe/H] in turn, and compare the
  cross-validated errors. Which representation spends more of its capacity on
  chemistry?
\end{enumerate}
\end{issues}
